\chapter{킬 조건부 교전 국소화}
\label{ch:localization}

이 장은 정의 \ref{def:engagement}를 계산하는 알고리즘(코드명 \texttt{teamfight\_v2})을 기술한다. 알고리즘은 킬 사건만을 씨앗으로 사용하며(kills-only creation principle), 와드·오브젝트 사건은 교전을 만들지 않는다. 전체 절차는 알고리즘 \ref{alg:localization}에 요약되어 있다.

\section{5 초 위치 격자}
\label{sec:loc-grid}

온셋 검증(제 \ref{sec:loc-validate} 절)은 온셋 시각 $t_e$에서 열 명의 위치를 요구하지만 프레임은 60 초 간격이다. 따라서 검출 전용의 5 초 격자를 만든다. 프레임 $F_k$와 $F_{k+1}$ 사이의 격자점 $t$에 대해
\begin{equation}
  \alpha_{\mathrm{raw}} = \frac{t - t_k}{t_{k+1} - t_k}, \qquad
  \tilde{\alpha} = 1 - e^{-\kappa\,\alpha_{\mathrm{raw}}}, \qquad
  \widehat{\mathrm{pos}}_i(t) = (1-\tilde\alpha)\,\mathrm{pos}_i(t_k) + \tilde\alpha\,\mathrm{pos}_i(t_{k+1})
  \label{eq:interp}
\end{equation}
로 위치를 보간한다($\kappa = 3$). 지수 곡선은 플레이어가 프레임 직후에 빠르게 다음 위치로 이동하는 경향을 반영한 경험적 선택이다. 여기에 \emph{킬 전 덮어쓰기}(pre-kill override)를 적용한다. 킬 $u$의 참가자 $i \in \mathrm{part}(u)$에 대해 직전 프레임 시각 $t_k$부터 $t(u)$까지의 구간에서
\begin{equation}
  \widehat{\mathrm{pos}}_i(t) = (1-\alpha_u)\,\mathrm{pos}_i(t_k) + \alpha_u\,\mathrm{pos}(u), \qquad \alpha_u = \frac{t - t_k}{t(u) - t_k}
\end{equation}
로 대체하여, 킬이 일어난 지점으로 참가자가 수렴하도록 만든다. 이 격자는 가정 \ref{asm:nokillpos}에 따라 \textbf{검출에만} 쓰이며 모델 입력이 되지 않는다.

\section{시간적 킬 군집화}
\label{sec:loc-temporal}

시각 순으로 정렬된 킬 열 $k_1, \dots, k_M$을 연속한 킬 사이의 간격이 $\delta_{\mathrm{gap}} = 18{,}000$ ms 이하이면 같은 군집으로 묶는다:
\begin{equation}
  k_j, k_{j+1} \in \text{같은 군집} \iff t(k_{j+1}) - t(k_j) \le \delta_{\mathrm{gap}}.
\end{equation}
이는 1차원 시간 축에서의 단일 연결(single-linkage) 군집화와 같으며, 한 번의 선형 순회로 계산된다. 18 초라는 값은 경쟁 LoL에서 전형적인 교전의 지속 시간과 부활 대기시간을 고려하여 정하였다.

\section{공간 지름 분할}
\label{sec:loc-spatial}

시간 군집은 서로 다른 장소에서 동시에 벌어진 교전(예: 탑 라인 소규모 교전과 봇 라인 다이브)을 하나로 묶을 수 있다. 이를 분리하기 위해 군집의 \emph{공간 지름}
\begin{equation}
  \mathrm{diam}(C) = \max_{u, u' \in C} \big\|\mathrm{pos}(u) - \mathrm{pos}(u')\big\|_2
\end{equation}
이 $D_{\max} = 4{,}000$ 단위를 넘으면 지름이 $D_{\max}$ 이하인 부분 군집들로 나눈다. 이는 완전 연결(complete-linkage) 기준의 응집적 군집화 \cite{sorensen1948method,defays1977efficient}에 해당한다. 단일 연결을 쓰면 킬이 사슬처럼 이어질 때 총 지름이 $D_{\max}$를 초과할 수 있으므로, 완전 연결이 정의 \ref{def:engagement}의 ``지름 상한''을 정확히 구현한다.\footnote{초기 구현은 단일 연결이었으며, 사전 감사에서 완전 연결로 수정되었다(제 \ref{sec:disc-audit} 절).} 군집의 모든 쌍 거리가 $D_{\max}$ 이하이면 분할 없이 통과하는 빠른 경로를 둔다.

\section{온셋 검증}
\label{sec:loc-validate}

각 후보 $C$에 대해 온셋 $t_e = t(k_1) - \delta_{\mathrm{pre}}$ ($\delta_{\mathrm{pre}} = 10{,}000$ ms, 경기 시작 시각으로 클램핑)를 계산하고 다음을 모두 요구한다.

\begin{enumerate}[leftmargin=2em]
  \item \textbf{경계 조건.} $t_e \ge t_{\min} + c_{\mathrm{ctx}}$이고 $t_e + h \le t_{\max}$ (관측 창과 라벨 창이 경기 안에 온전히 들어감).
  \item \textbf{생존 조건.} $t_e$에서 양 팀 각각 생존 챔피언이 $n_{\min} = 2$명 이상.
  \item \textbf{공간 조건.} 중심 $c = \mathrm{pos}(k_1)$에 대해
  \begin{equation}
    \big|\{i \in \mathrm{Blue} : \mathrm{alive}_i(t_e) \wedge \|\widehat{\mathrm{pos}}_i(t_e) - c\| \le r_{\mathrm{val}}\}\big| \ge n_{\min}, \quad
    \big|\{i \in \mathrm{Red} : \cdots\}\big| \ge n_{\min}.
    \label{eq:validate}
  \end{equation}
\end{enumerate}
식 \eqref{eq:validate}는 ``생존''과 ``반경 안''을 \emph{하나의 논리곱}으로 검사한다. 즉 사망한 채 반경 안에 남아 있는 챔피언은 계수되지 않는다. 이 조건은 한 명이 잡히는 픽(pick)이나 일방적인 정리(cleanup)를 교전에서 배제한다.

\section{인접 후보 병합}
\label{sec:loc-merge}

검증을 통과한 후보 $C_a, C_b$가 시간적으로 $\delta_{\mathrm{merge}} = 15{,}000$ ms 이내이고 중심 거리가 $r_{\mathrm{merge}} = 2{,}000$ 단위 이내이면 하나의 교전으로 병합한다. 병합 후 길이가 $T_{\max} = 60{,}000$ ms를 넘는 교전은 기각한다. 이 단계는 한 교전 안에서 킬 사이의 간격이 잠깐 18 초를 넘었다가 같은 장소에서 다시 이어지는 경우(재교전)의 중복 계수를 막는다.

\section{교전 유형 분류와 부가 정보}
\label{sec:loc-classify}

병합된 교전에는 두 축의 유형이 부여된다. \emph{규모}(fight scale)는 근접 쌍의 수로 정의하여 8 쌍 이상이면 한타(teamfight), 4 쌍 이상이면 소규모 교전(skirmish), 그 미만이면 픽(pick)으로 하고, \emph{맥락}(fight context)은 중심 위치가 바론·용·전령 둥지 1{,}500 단위 이내, 포탑 1{,}000 단위 이내, 본진 3{,}000 단위 이내인지에 따라 오브젝트 교전, 포탑 다이브, 본진 교전 등으로 나눈다. 또한 교전 중 반경 3{,}000 단위 안의 비-킬 상호작용(와드, 소환사 주문)과, 마지막 킬 이후 30 초 동안의 오브젝트·포탑·골드 변동을 부가 정보로 기록한다. 이들은 하위 집단 분석에만 쓰이며 라벨이나 입력이 아니다.

\section{알고리즘과 복잡도}
\label{sec:loc-algorithm}

\begin{algorithm}[htbp]
\caption{킬 조건부 교전 국소화 (\texttt{teamfight\_v2})}
\label{alg:localization}
\begin{algorithmic}[1]
\Require 프레임 열 $\mathcal{F}_m$, 사건 열 $\mathcal{E}_m$, 상수 $\delta_{\mathrm{gap}}, D_{\max}, \delta_{\mathrm{pre}}, r_{\mathrm{val}}, n_{\min}, \delta_{\mathrm{merge}}, r_{\mathrm{merge}}, T_{\max}$
\Ensure 유효 교전 집합 $\mathcal{G}_m$
\State $\widehat{\mathrm{pos}} \gets$ 5초 위치 격자 (식 \eqref{eq:interp}) + 킬 전 덮어쓰기 \Comment{검출 전용}
\State $\mathcal{K} \gets$ 시각 순으로 정렬된 \texttt{CHAMPION\_KILL} 사건
\State $\mathcal{C} \gets$ 연속 간격 $\le \delta_{\mathrm{gap}}$인 킬들의 시간 군집 \Comment{단일 연결, $O(M)$}
\ForAll{$C \in \mathcal{C}$ with $\mathrm{diam}(C) > D_{\max}$}
  \State $C$를 완전 연결 기준으로 지름 $\le D_{\max}$인 부분 군집으로 분할 \Comment{$O(|C|^2)$}
\EndFor
\State $\mathcal{V} \gets \emptyset$
\ForAll{후보 $C$}
  \State $t_e \gets \max(t(k_1) - \delta_{\mathrm{pre}},\, t_{\min})$;\quad $c \gets \mathrm{pos}(k_1)$
  \If{경계 조건 $\wedge$ 생존 조건 $\wedge$ 식 \eqref{eq:validate}}
    \State $\mathcal{V} \gets \mathcal{V} \cup \{(t_e, t(k_1), t(k_{|C|}), c, C)\}$
  \EndIf
\EndFor
\State $\mathcal{G}_m \gets$ $\mathcal{V}$에서 시간 $\le \delta_{\mathrm{merge}}$, 거리 $\le r_{\mathrm{merge}}$인 인접 후보를 병합; 길이 $> T_{\max}$ 기각
\State \Return $\mathcal{G}_m$
\end{algorithmic}
\end{algorithm}

한 경기의 킬 수를 $M$, 프레임 수를 $T$라 하면 격자 구축은 $O(NT \cdot 12)$, 시간 군집화는 $O(M)$, 지름 분할은 군집 크기 $|C|$에 대해 $O(|C|^2)$이므로 최악 $O(M^2)$, 검증은 후보당 $O(N)$, 병합은 $O(|\mathcal{V}|)$이다. 경기당 킬 수는 수십 건 수준이므로 전체는 사실상 프레임 처리에 지배되며, 206{,}442 경기 전체를 단일 워크스테이션에서 처리할 수 있다. 알고리즘은 결정론적이므로 같은 입력에서 항상 같은 교전 집합을 낸다.

\section{검출기 타당성 검증 설계}
\label{sec:loc-validation}

알고리즘 \ref{alg:localization}은 사람의 주석 없이 말뭉치를 만들지만, 그 결과가 사람이 ``교전''이라 부르는 것과 일치하는지는 별도로 검증해야 한다(RQ1). 본 논문은 두 트랙의 검증을 설계하였다.

\paragraph{트랙 A: 텔레메트리 재생 기반 주석 연구.} 검출기가 소비하는 것과 동일한 Match-V5 텔레메트리로 미니맵 재생을 재구성하고, 주석자가 검출기 출력을 보지 않은 상태(blind)에서 모든 교전 구간을 표시한다. 이 트랙은 알고리즘의 \emph{결정 규칙}(18 초, 4{,}000 단위, 1{,}800 단위, 15 초/2{,}000 단위)이 사람의 판단과 일치하는지를 검증한다. 다만 API에 없는 정보(킬 없는 대치)는 검출할 수 없으므로 별도 항목(\texttt{killless\_annotated})으로 보고한다.

\paragraph{트랙 B: 전체 리플레이 골드셋.} 2026년 KR 랭크 경기(공개 패치 26.13, API 패치 16.13) 500 건의 완전한 리플레이(\texttt{.rofl})를 수집하고, 리플레이·detail·timeline 세 파일의 SHA-256을 대조하였다. 이 가운데 10 분 미만 경기 12 건을 제외하고, 생성 시각 10 개 구간에서 각 10 경기씩 결정론적으로(시드 20260715) 100 경기의 골드셋을 뽑았다. 표본 추출에 검출기 출력은 사용하지 않았다. 검출기는 이 100 경기에서 480 개의 후보 구간(경기당 4.8 개)을 산출하였고, 실행 예외·반복 실행 불일치·경계 오류는 없었다.

두 트랙 모두 두 명의 주석자가 독립적으로 교전의 \emph{시작}(첫 킬이 아닌 교전 개시)부터 \emph{종료}(이탈 또는 전멸)까지를 표시하며, 결과는 시간 IoU 0.5 기준의 탐욕적 일대일 짝짓기로 정밀도·재현율·$F_1$·평균 IoU·평균 온셋 오차를 계산하고, 주석자 간 $F_1$과 Cohen의 $\kappa$를 함께 보고한다. 트랙 B의 사람 주석은 본 논문 작성 시점에 진행 중이며(제 \ref{sec:results-detector} 절), 기술적 통과 자체를 검출 정확도로 인용하지 않는다.
